\documentclass[10pt]{article}
\usepackage[T1]{fontenc}
\usepackage[utf8]{inputenc}
\usepackage[french]{babel}
\usepackage{amsmath,amssymb,mathrsfs,mathpazo}
\usepackage{geometry}
\geometry{a4paper,margin=2cm}
\usepackage{graphicx,array,tabularx,tabularray,colortbl,multirow,multicol,booktabs,enumitem,pifont,ulem,fancybox,diagbox,cancel}
\usepackage{pstricks,pstricks-add,pst-plot,pst-tree,pst-node,pst-text,pst-func,pst-math,pst-bspline,pst-3dplot,pst-eucl,pst-solides3d}
\let\pscalculate\undefined
\usepackage{pst-bezier}
\usepackage{tikz,pgfplots,tkz-tab}
\usetikzlibrary{arrows,arrows.meta,patterns,calc,positioning,trees,decorations.markings,shapes.geometric,shapes.misc,shapes.arrows,intersections,decorations.pathmorphing,decorations.pathreplacing,backgrounds,shadows,babel}
\pgfplotsset{compat=1.18}
\usepgfplotslibrary{fillbetween}
\usepackage[np]{numprint}
\usepackage{listings,makecell,pas-tableur,twoopt}
\let\barre\undefined
\usepackage[pstricks]{tablvar}
\definecolor{darkgreen}{RGB}{0,100,0}
\providecommand{\red}{\color{red}}
\providecommand{\blue}{\color{blue}}
\providecommand{\green}{\color{green}}
\providecommand{\black}{\color{black}}
\providecommand{\white}{\color{white}}
\providecommand{\gray}{\color{gray}}
\providecommand{\cyan}{\color{cyan}}
\providecommand{\magenta}{\color{magenta}}
\providecommand{\yellow}{\color{yellow}}
\providecommand{\violet}{\color{violet}}
\providecommand{\orange}{\color{orange}}
\providecommand{\euro}{\texteuro}
\providecommand{\EUR}{\texteuro}
\providecommand{\ieme}{\textsuperscript{e}}
\providecommand{\degres}{\ensuremath{^\circ}}
\providecommand{\degree}{\ensuremath{^\circ}}
\providecommand{\vv}[1]{\overrightarrow{#1}}
\providecommand{\Coord}[1]{#1}
\providecommand{\decosix}{}
\providecommand{\starredbullet}{\ensuremath{\bullet}}
\pagestyle{empty}
\setlength{\parindent}{0pt}
\usepackage{adjustbox}\DeclareUnicodeCharacter{2236}{:}\DeclareUnicodeCharacter{2190}{\ensuremath{\leftarrow}}
\providecommand{\R}{\mathbb{R}}
\providecommand{\N}{\mathbb{N}}
\providecommand{\Z}{\mathbb{Z}}
\providecommand{\Q}{\mathbb{Q}}
\providecommand{\D}{\mathbb{D}}
\providecommand{\C}{\mathbb{C}}
\providecommand{\e}{{\mathrm{e}}}
\providecommand{\vect}[1]{\overrightarrow{#1}}
\providecommand{\vectt}[1]{\overrightarrow{\mathrm{#1}}}
\providecommand{\vvt}[1]{\overrightarrow{\mathrm{#1}}}
\providecommand{\np}[2][]{{#2~#1}}
\providecommand{\vv}[1]{{\overrightarrow{#1}}}
\renewcommand{\d}{{\mathrm{d}}}
\providecommand{\degres}{{^\circ}}
\providecommand{\degre}{{^\circ}}
\providecommand{\blue}{\color{blue}}
\providecommand{\red}{\color{red}}
\providecommand{\enskip}{\quad}
\providecommand{\llbracket}{\lbrack\!\lbrack}
\providecommand{\rrbracket}{\rbrack\!\rbrack}
\providecommand{\euro}{{\text{€}}}
\providecommand{\no}{{\text{n°}}}
\providecommand{\interval}[2]{{[#1;#2]}}
\providecommand{\intervFO}[2]{{[#1;#2[}}
\providecommand{\dotfill}{\ldots}

\providecommand{\ds}{\displaystyle}
\providecommand{\textding}[1]{#1}
\providecommand{\Oij}{\(\mathrm{O};\vec{\imath},\vec{\jmath}\)}
\providecommand{\Oijk}{\(\mathrm{O};\vec{\imath},\vec{\jmath},\vec{k}\)}
\providecommand{\Ouv}{\(\mathrm{O};\vec{u},\vec{v}\)}
\providecommand{\intervFF}[2]{\ensuremath{[#1;#2]}}
\renewcommand{\euro}{\ifmmode\text{\texteuro}\else\texteuro\fi}
\DeclareUnicodeCharacter{00B0}{\ifmmode^{\circ}\else\textdegree\fi}
\DeclareUnicodeCharacter{20AC}{\ifmmode\text{\texteuro}\else\texteuro\fi}
% Énoncés : Éducation nationale. Transcription et corrigés : APMEP, Denis Vergès et contributeurs. Découpage et classement : CoopMaths, Sébastien Lozano, Éric Elter et contributeurs. Adaptation de la présentation : Maths Explorer.
% CC BY-SA : https://coopmaths.fr/www/about/
% https://coopmaths.fr/alea/static/bac/2023/tex/bac_2023_03_sujet1_ameriquenord_4_cor.tex
\begin{document}
\noindent\textit{APMEP / CoopMaths. CC BY-SA. Présentation adaptée par Maths Explorer.}\par\medskip
\medskip

\textbf{Partie A - Étude de la suite \boldmath{$\left(u_{n}\right)$}}

\medskip

\begin{enumerate}
\item Calculons $u_{1}$ :\quad $u_1 = \dfrac{1}{2}\left(u_0 + \dfrac{11}{u_0}\right) = \dfrac{1}{2}\left(5 + \dfrac{11}{5}\right) = \dfrac{1}{2}\times \dfrac{25 + 11}{5} = \dfrac{18}{5}$.

\hphantom{Calculons $u_{1}$ :\quad} $u_2 = \dfrac{1}{2}\left(u_1 + \dfrac{11}{u_1}\right) = \dfrac{1}{2}\left(\dfrac{18}{5} + \dfrac{11}{\frac{18}{5}}\right) = \dfrac{1}{2}\times \left(\dfrac{18}{5}+\dfrac{55}{18}\right)  = \dfrac{599}{180}$.

\item La fonction $f$ est dérivable sur son ensemble de définition (c'est une fraction rationnelle).

Pour tout réel strictement positif, on a:

$f'(x) = \dfrac{1}{2}\left(1 + 11\times \dfrac{-1}{x^2}\right) = \dfrac{1}{2}\times \dfrac{x^2 - 11}{x^2} = \dfrac{x^2-11}{2x^2}$

Pour tout $x$ dans l'intervalle $\left[\sqrt{11} ~;~+\infty\right[$, on a :

$\begin{aligned} x \geqslant \sqrt{11} &\implies x^2 \geqslant 11 \qquad\text{ car la fonction carré est croissante sur }\R^+\\
&\implies x^2 - 11 \geqslant 0\\
&\implies \dfrac{x^2-11}{2x^2} \geqslant 0 \qquad \text{car, pour tout réel non nul, } 2x^2>0 \\
&\implies f'(x) \geqslant 0\\ \end{aligned}$

Sur l'intervalle $\left[\sqrt{11} ~;~+\infty\right[$, $f'$ est à valeurs positives, donc $f$ est croissante.

\item \emph{Remarque :} on constate que la fonction $f$ est telle que, pour tout entier naturel $n$, on a :

$u_{n+1} = f(u_n)$, et par ailleurs, on va calculer $f\left(\sqrt{11}\right) = \dfrac{1}{2}\left(\sqrt{11} + \dfrac{11}{\sqrt{11}}\right)  = \dfrac{1}{2}\left(\sqrt{11} + \sqrt{11}\right) = \sqrt{11}$

Pour tout entier naturel $n$, on s'intéresse à l'inégalité : $ u_{n} \geqslant u_{n+1} \geqslant \sqrt{11}$.

\emph{Initialisation :} Pour $n = 0$, on a $u_n = 5$ et $u_{n+1} = u_1 = \dfrac{18}{5} = 3,6$ et par ailleurs, on a $\sqrt{11}\approx 3,32$ donc on a bien $u_0\geqslant u_1 \geqslant\sqrt{11}$. L'inégalité est vraie à l'indice 0.

		\emph{Hérédité :} Pour un entier naturel $n$ quelconque fixé, on suppose que l'inégalité est vraie pour l'indice $n$, c'est-à-dire que l'on a : $ u_{n} \geqslant u_{n+1} \geqslant \sqrt{11}$.

		 $\begin{aligned} u_{n} \geqslant u_{n+1} \geqslant \sqrt{11}
		 	&\implies f(u_n) \geqslant f(u_{n+1}) \geqslant f\left(\sqrt{11}\right) \qquad \text{car  }f\text{ est croissante sur }\left[\sqrt{11}~;~+\infty \right[\\
		 	&\implies u_{n+1}\geqslant u_{n+2} \geqslant \sqrt{11} \qquad \text{grâce aux éléments de la remarque}\\ \end{aligned}$

\emph{Conclusion :} On a montré que l'inégalité est vraie pour l'indice 0, et que, pour tout entier naturel $n$ la véracité est héréditaire de l'indice $n$ à l'indice $(n+1)$. En vertu du principe de récurrence, on peut en déduire que :

Pour tout $n$ naturel, on a : $u_{n} \geqslant u_{n+1} \geqslant \sqrt{11}$.

\item On a notamment :
\begin{itemize}
\item $\forall n \in \N, \quad u_n\geqslant \sqrt{11}$ : la suite $(u_n)$ est minorée par $\sqrt{11}$.

\item $\forall n \in \N, \quad u_n\geqslant u_{n+1}$ : la suite $(u_n)$ est décroissante.
\end{itemize}

D'après le théorème de convergence monotone, la suite est donc convergente vers une limite $a$, telle que $a \geqslant\sqrt{11}$.

	\item $u$ est une suite convergente vers une limite réelle $a$ et $u$ a une relation de récurrence de la forme :

$\forall n \in \N, \quad u_{n+1}=f(u_n)$

où $f$ est une fonction continue sur $\R^{*+}$, ensemble qui contient tous les termes de la suite (car celle-ci est minorée par $\sqrt{11}>0$).

D'après le théorème du point fixe, on en déduit que la limite $a$ de la suite doit être une solution de l'équation $f(x) = x$.

Résolvons cette équation dans $\R^{*+}$ :

$\begin{aligned}  f(x) = x &\iff \dfrac{1}{2}\left(x + \dfrac{11}{x}\right) = x\\
&\iff x+ \dfrac{11}{x} = 2x\\
&\iff -x + \dfrac{11}{x} = 0\\
&\iff -x^2 + 11 = 0 \qquad \text{car sur }\R^{*+}, \quad x>0\\
&\iff x^2 = 11\\
&\iff x = \sqrt{11} \qquad \text{car on résout sur }\R^{*+}		\\ \end{aligned}$

L'équation n'a qu'une unique solution sur $\R^{*+}$, c'est $\sqrt{11}$

On en conclut donc que la suite ne peut converger que vers $\sqrt{11}$.

On a donc $\lim\limits_{n \to +\infty} u_n = \sqrt{11}$.
\end{enumerate}

\medskip

\textbf{Partie B - Application géométrique}

\medskip

\begin{enumerate}
\item 
	\begin{enumerate}
		\item Le rectangle $R_0$ a pour aire 11 et pour longueur $L_0 = 5$, donc sa largeur $\ell_{0}$ est telle que $L_0 \times \ell_0 = 11 \iff \ell_0 = \dfrac{11}{L_0}$

On a donc bien $\ell_{0}=\dfrac{11}{5} = 2,2$.

		\item De façon général, pour tout $n$ entier naturel :

		Le rectangle $R_n$ a pour aire 11 et pour longueur $L_n$, donc sa largeur $\ell_{n}$ est telle que $L_n \times \ell_n = 11 \iff \ell_n = \dfrac{11}{L_n}$.

	\end{enumerate}
\item On sait que pour tout $n$ naturel : $L_{n+1} = \dfrac{L_n + \ell_n}{2} = \dfrac{1}{2}\left(L_n + \ell_n\right)  = \dfrac{1}{2}\left(L_n + \dfrac{11}{L_n}\right)$.

La relation de récurrence vérifiée par la suite $(L_n)$ est donc la même que celle vérifiée par la suite $(u_n)$ dans la \textbf{partie A}. Comme de plus $L_0 = u_0 = 5$, les deux suites sont donc rigoureusement égales.
	\item À la \textbf{partie A}, question \textbf{3.}, on a établi que la suite $(u_n)$, donc la suite $(L_n)$ est minorée par $\sqrt{11}$.

On a donc $\sqrt{11}\leqslant L_n$. De plus :

$\begin{aligned}[t] \sqrt{11}\leqslant L_n &\implies \dfrac{1}{\sqrt{11}} \geqslant \dfrac{1}{L_n} \qquad\text{car la fonction inverse est strictement décroissante sur }\R^{*+}\\
&\implies \dfrac{11}{\sqrt{11}} \geqslant \dfrac{11}{L_n} \qquad \text{car } 11>0\\
&\implies \sqrt{11}\geqslant\ell_n\\
\\ \end{aligned}$

En prenant cette inégalité et celle dont on est parti, on a bien : $\ell_n \leqslant\sqrt{11} \leqslant L_n$.

\item On admet que les suites $\left(L_{n}\right)$ et $\left(\ell_{n}\right)$ convergent toutes les deux vers $\sqrt{11}$.

Géométriquement, cela veut dire que notre rectangle $R_n$ a des dimensions qui tendent vers $\sqrt{11}$, donc le rectangle tend vers un carré de côté $\sqrt{11}$.

item Ce script renvoie un encadrement de $\sqrt{11}$ par les valeurs $\ell_n$ et $L_n$, arrondies au millionième près, où $n$ est l'argument d'appel de la fonction.

	\begin{enumerate}
		\item Pour l'appel \texttt{heron(3)} le retour de la fonction sera donc $\ell_3$ et $L_3$, arrondis au millionième près.

on obtient donc les valeurs suivantes : \np{3,316606} , \np{3,316643}

		\item Une interprétation de ces deux valeurs, c'est donc que le nombre $\sqrt{11}$ est compris entre \np{3,316606} et \np{3,316643}, soit un encadrement d'un amplitude inférieure à $4\times 10^{-5}$, soit une précision importante, avec seulement trois itérations de notre algorithme.
		\end{enumerate}

	\end{enumerate}
\emph{Remarque :} Ici, la fonction python a été appelée \texttt{heron} en référence à Héron d'Alexandrie (qui a vécu sans doute au cours du Ier siècle après JC.).

La méthode de Héron ou méthode babylonienne est une méthode efficace d'extraction de racine carrée, c’est-à-dire de résolution de l'équation $x^2 = a$, avec $a$ positif. Elle porte le nom du mathématicien Héron d'Alexandrie, qui l'expose dans le tome I de son ouvrage Metrica (Les métriques), découvert seulement en 1896 mais certains calculs antérieurs, notamment égyptiens, semblent prouver que la méthode est plus ancienne.
\hfill(source : wikipedia)

\end{document}
